\begin{frame}{발산의 모습과 실전 대응}
  \begin{itemize}
    \item 비용이 줄어들기는커녕 스텝마다 \textbf{자릿수가 계속 늘어나며
      폭발} --- 걸음이 너무 커서 골짜기 바닥을 지나쳐 반대편 비탈로
      튕겨 나가고, 다음 스텝엔 그 반대편에서 \textbf{더 크게} 튕겨
      나오는 반복.
  \end{itemize}
  \vskip0.6em
  \begin{block}{실전 대응}
    손실이 \texttt{NaN} 이나 매우 큰 값으로 튀면, \textbf{가장 먼저
    의심할 것이 바로 학습률}.
    \\ 보통 $\alpha = 0.1, 0.01, 0.001$ 처럼 10배씩 줄여가며 손실
    곡선이 매끄럽게 내려가는 지점을 찾는다 --- 이 ``10배씩''
    경험칙이 임의가 아니라 데이터 특이적 경계값 $\alpha^{*}$ 가
    존재하기 때문 (앞 슬라이드 확인).
  \end{block}
\end{frame}
